% =====================================================================
%  13 -- BEYOND THE MAP: THE BOUNDARY AUDIT, AND THE EMULATOR AS AN OBJECT
%
%  Source: tier3.md lines 5204-5912, in full and in order --
%          Part IX header + intro (5204-5210)      -> sec:boundaryaudit
%          IX.1 (5212-5237)                        -> sec:swept
%          IX.2 (5239-5262)                        -> sec:noareas
%          IX.3 (5264)                             -> sec:threedeveloped
%            IX.3.1 (5266-5303)                    -> sec:multiplicity   (\paragraph)
%            IX.3.2 (5305-5348)                    -> sec:artifact       (\paragraph)
%            IX.3.3 (5350-5373)                    -> sec:nucleardata    (\paragraph)
%          IX.4 (5375-5400)                        -> sec:axescovered
%          IX.5 (5402-5426)                        -> sec:converging
%          Part X header + intro (5429-5439)       -> sec:emulatorobject
%          X.1 (5440-5522) .. X.11 (5860-5912)     -> sec:fixedpoints .. sec:convergingnow
%
%  Cross-reference rewordings (each also marked inline with % [edition]):
%    - IX intro "the reason several sections of Part 0 exist"
%        -> \S\ref{part:engine}--\S\ref{part:host}
%    - IX.5 "the current register's cheapest and highest-ranked items (T-24,
%      T-03, T-08)" -> "the register's cheapest and highest-ranked items
%      (\S\ref{sec:ranked}; \reg{T-24}, ...)" -- the register FOLLOWS this
%      section in the report.
%    - X intro "SSIX.5" -> \S\ref{sec:converging}; "SSX.11" -> \S\ref{sec:convergingnow}
%    - X.11 "(SSVII.2)" -> \S\ref{sec:retired} (forward reference; wording kept)
%    - X.11 "Pass 3 (Part IX)" -> \S\ref{sec:boundaryaudit}
%  All other section pointers into tier3.md are live \S\ref{...} per the
%  label registry; pointers into Tier 0/1/2 stay as plain text.
%  Blockquotes: 5623-5626 result; 5789-5794 aside; 5831-5834 plain quote
%  (from docs/reports/consolidated-gnn-architecture.md); 5883-5887 result.
% =====================================================================

\section{Beyond the map: the boundary audit, and the emulator as an object}
\label{part:beyond}

% =====================================================================
\subsection{Prerequisites outside the map: a cross-tier audit}
\label{sec:boundaryaudit}

The tier map (S0--S16, plus the sketched S17--S22) organises the project by
\emph{where the physics lives}. This part asks the orthogonal question:
\textbf{what is prerequisite to this project that the map has no node for at
all?} It is a second-pass audit across all four written tiers, and it is the
reason several sections of \S\ref{part:engine}--\S\ref{part:host} exist.
% [edition] reworded cross-ref: "the reason several sections of Part 0 exist"

% ---------------------------------------------------------------------
\subsubsection{What was swept, and how}
\label{sec:swept}

Three passes over different axes, deliberately chosen to be ones the map does not
organise by:

\begin{itemize}
\item \textbf{Pass A --- the physical setting outside the burning zone.}
  Everything a stellar-evolution model does that is \emph{not} the network:
  transport, structure, the progenitor's history, and the collapse's numerics.
  Produced \reg{T-19} (mixing), \reg{T-21} (the NSE transition), \reg{T-22}
  (the progenitor grid), \reg{T-23} (what actually runs these calculations).
\item \textbf{Pass B --- the interface.} What consumes the emulator's output, in
  the concrete sense of a function signature. Produced \reg{T-24} (the
  bottleneck fraction and Amdahl), \reg{T-25} (host-side derivatives),
  \reg{T-26} (the lost convergence signal), \reg{T-27} (per-call latency and
  interop).
\item \textbf{Pass C --- the mathematics of the answer, and the practice of
  producing it.} The geometry of the output space, the statistics of the
  claims, and the lifecycle of the artifact. Produced \reg{T-28} (float32 vs
  cancellation), \reg{T-30} (the simplex), \reg{T-31} (multiplicity),
  \reg{T-32} (the artifact lifecycle), plus \reg{T-20} (the error hierarchy)
  and \reg{T-29} (the fixed-point audit) as methodological items.
\end{itemize}

Each pass asked the same three questions of its axis: \emph{what does the
project assume here? what would a practitioner in this area immediately ask? and
is there any test that could currently fail?}

% ---------------------------------------------------------------------
\subsubsection{The areas with no study node anywhere}
\label{sec:noareas}

\begingroup
\footnotesize
\begin{longtable}{@{}L{2.9cm} L{4.6cm} L{4.0cm} L{2.6cm}@{}}
\toprule
\textbf{area} & \textbf{prerequisite because} & \textbf{nearest existing coverage} & \textbf{now at} \\
\midrule
\endfirsthead
\toprule
\textbf{area} & \textbf{prerequisite because} & \textbf{nearest existing coverage} & \textbf{now at} \\
\midrule
\endhead
\bottomrule
\endfoot

\textbf{Convection and mixing} &
  the network's output is immediately mixed; the burning/mixing competition
  decides whether \Ye{} errors average out &
  Tier~0 \S 0.7's timescale row + a one-line caveat &
  \S\ref{sec:convection}, \reg{T-19} \\
\textbf{The NSE transition} &
  the host replaces the network with a table exactly where the labels are
  pathological &
  \tierone{VIII.E.6} (``the handoff is unspecified'') &
  \S\ref{sec:nsetransition}, \reg{T-21} \\
\textbf{Progenitor diversity} &
  every real-track statement rests on one 20\,\Msun{} model &
  Tier~2 \S V.2 (that one model) &
  \S\ref{sec:whichstars}, \reg{T-22} \\
\textbf{What computes explodability} &
  \reg{T-01} asks for a derivative of it &
  nothing &
  \S\ref{sec:explodability}, \reg{T-23} \\
\textbf{The error hierarchy} &
  decides whether the improvement is scientifically visible &
  scattered; never assembled &
  \S\ref{sec:errorhierarchy}, \reg{T-20} \\
\textbf{The runtime fraction $p$} &
  the project's founding premise; Amdahl caps everything &
  Tier~0 \S I.3 derives the \emph{scaling}, never the fraction &
  \S\ref{sec:bottleneck}, \reg{T-24} \\
\textbf{The call-site contract} (7 items) &
  an emulator that cannot be installed is not a result &
  \foreignreg{2}{R-11} covers 2 of the 7 &
  \S\ref{sec:callsite}, \S\ref{sec:netreturns}, \reg{T-25}/\reg{T-26}/\reg{T-27} \\
\textbf{Simplex geometry / compositional data} &
  it is the space the output lives in, with a mature literature &
  Tier~2 \S I.7b (positivity, as a physics worry) &
  \S\ref{sec:simplex}, \reg{T-30} \\
\textbf{Experiment design \& multiplicity} &
  18 cells $\times$ 6 metrics makes ``we win'' cheap &
  nothing &
  \S\ref{sec:multiplicity}, \reg{T-31} \\
\textbf{Artifact lifecycle} &
  what ships, and whether it reproduces &
  \foreignreg{2}{R-25} (shelf life, physics side only) &
  \S\ref{sec:artifact}, \reg{T-32} \\
\end{longtable}
\endgroup

Ten areas. Note the shape: \textbf{one is physics-internal (mixing), four are
about the physical setting the project sits inside, three are about the
interface, and two are about the practice of producing and shipping the result.}
The map's blind spot is not the physics --- Tiers 0--2 are thorough --- it is
everything on the boundary of the physics.

% ---------------------------------------------------------------------
\subsubsection{Three developed here, having no home section}
\label{sec:threedeveloped}

\paragraph{Experiment design and the multiplicity problem}
\label{sec:multiplicity}

The project will compare an emulator against a baseline across
\textbf{2 networks $\times$ 9 $\Delta t$ $\times$ 6 metrics = 108 cells}, and
it has already done exactly that once, for the reproduction
(\resultsref{2026-07-08}, ``ratio 1.000 in all 108 comparisons''). For a
\emph{reproduction} that is a strength --- 108 exact matches is overwhelming
evidence. For a \emph{comparison} it is a trap.

\textbf{The arithmetic.} If two models are genuinely equivalent and each cell is
a coin flip, the probability that one wins \textbf{at least 60 of 108} cells is
about 15\% (at least 63 of 108: about 5\% --- corrected in the 2026-08-18 audit;
the earlier text put 5\% at 60), and the probability of winning at least one
cell ``by a lot'' is essentially 1. With per-cell noise from training seeds and
test-set sampling, a table of 108 numbers will always contain a flattering
subset. \derivedhere[binomial]

\textbf{What the project needs and does not have:}

\begin{enumerate}
\item \textbf{A pre-registered primary metric.} One number, chosen before
  training, against which the claim is made. Everything else is secondary and
  labelled as such. The obvious candidate is absolute $|\dYe|$ p99 on the
  relaxed manifold in the QSE window --- it is the project's own quantity, on
  the deployment distribution, in the regime it claims.
\item \textbf{Seed replication.} \code{CLAUDE.md} already requires three seeds
  for the \emph{learned-mask regression} trigger, which shows the instinct is
  present; it is not required for the headline comparison. Without $\ge 3$
  seeds there is no way to distinguish a model improvement from an
  initialisation.
\item \textbf{Correction for multiplicity} on the secondary table --- or, more
  honestly, the discipline of reporting the whole table and claiming nothing
  from individual cells.
\item \textbf{A pre-declared definition of ``beating''}
  (\S\ref{sec:protocol}'s open decision 3), because deciding it after seeing the
  numbers is how a null result becomes a paper.
\end{enumerate}

None of this is exotic; it is standard experimental practice imported from
fields that had a replication crisis. The project's \code{RESULTS.md} discipline
--- every number dated, scripted, and provenanced --- is \emph{exactly} the
infrastructure that makes pre-registration cheap, and it is one line of policy
away from having it. Registered as \reg{T-31}.

\paragraph{The artifact lifecycle --- what actually ships}
\label{sec:artifact}

\foreignreg{2}{R-25} asks what a rate-library or MESA version bump does to a
trained emulator. That is the \emph{physics} half of a larger question the
project has not posed: \textbf{what is the deliverable, and can someone else run
it in five years?}

Six components, with the project's current status:

\begingroup
\small
\begin{tabularx}{\textwidth}{@{}Y L{6.4cm}@{}}
\toprule
\textbf{component} & \textbf{status} \\
\midrule
the trained weights & n/a (pre-training) \\
the \textbf{exact $\nu$, $C$, isotope order, and column order} the weights were
  trained against &
  \textbf{excellent} --- exported deterministically with content hashes,
  \code{disposition\_sha256} in the npz \\
the rate configuration (REACLIB snapshot, weaklib ordering, screening) &
  \textbf{excellent} --- pinned, reconciled, measured \\
a \textbf{serving artifact} with no Python in the loop &
  \textbf{nonexistent} (\S\ref{sec:callsite} item~5, \reg{T-27}) \\
\textbf{determinism} across thread counts, devices, and batch sizes &
  \textbf{unspecified} (\tierone{VIII.E.4}) \\
\textbf{model identity recorded in the consumer's output} &
  \textbf{unspecified} (\S\ref{sec:callsite} item~7) \\
\bottomrule
\end{tabularx}
\endgroup

The first three are as good as any project of this kind. The last three are
absent, and they are the ones that decide whether the work is \emph{usable}
rather than merely \emph{correct}.

Two specific traps worth naming because they are cheap to avoid in advance and
expensive to fix later:

\begin{itemize}
\item \textbf{Batch-size-dependent output.} Many frameworks produce
  bitwise-different results at different batch sizes (different GEMM kernels,
  different reduction orders) \provtag{[engineering, assumed --- no source
  opened]}. A stellar code that batches a variable number of zones per timestep
  will therefore get different answers for the same zone depending on how many
  neighbours were in the batch --- and, worse, a \emph{retry}
  (\S\ref{sec:callsite} item~4) may batch differently. This must be a declared,
  tested contract, not a hope.
\item \textbf{The conservation guarantee must survive export.}
  $\dd Y = \nu\netflux$ in float64 is exact; the same computation in a float32
  inference graph is not. The conservation map has to be pinned to float64
  \emph{in the exported artifact}, which is a concrete constraint on how the
  model is compiled, and \code{CLAUDE.md}'s ``float32 in model internals only''
  rule is the policy that anticipates it --- but nothing tests it end to end.
\end{itemize}

Registered as \reg{T-32}. It is an ADR plus a test, not a research programme.

\paragraph{Nuclear-data uncertainty is probably the binding term --- promote it}
\label{sec:nucleardata}

This is not a new area; it is \tierone{VIII.C.1} (rate uncertainty $\to$ \Ye{}
propagation) and \S VIII.C.5 (mass/$Q$ provenance). The audit's contribution is
to argue that they should be \textbf{first}, not twelfth.

The argument, in three steps:

\begin{enumerate}
\item \S\ref{sec:errorhierarchy}'s hierarchy puts the emulator's target
  ($3\times10^{-6}$/step) \emph{below} the weak-library shift
  ($5\times10^{-3}$/trajectory) by design. That is the intended relationship ---
  but it is a comparison against \textbf{one} known systematic.
\item The \emph{unknown} systematics are the Hauser--Feshbach rate uncertainties
  on unmeasured channels, which Tier~1 \S 1.3.3 puts at factor-of-two scale for
  the experimentally inaccessible ones (\citealt{rauscher2000} estimate their
  rates ``within a factor of 1.5--2''), and the mass/$Q$ differences, which
  Tier~1 \S VIII.C.5 notes enter as $e^{-Q/kT}$ with $Q/kT \approx 20$--50 for
  individual channels ($Q \approx 7$--13\,MeV at $kT \approx 0.26$--0.43\,MeV;
  hundreds to ${\sim}2000$ only for total binding energies in the Saha exponent)
  --- a 100\,keV $Q$ difference is a \textbf{47\% rate difference} at
  $kT = 0.26$\,MeV.
\item Nobody has propagated either to \Ye. If the propagated uncertainty is
  $10^{-3}$ per trajectory, the gate is well placed. If it is $10^{-2}$, the
  gate is over-engineered by orders and \textbf{the entire cost wall of Tier~2
  \S IV.9 is being paid for nothing.}
\end{enumerate}

\textbf{Everything needed to answer it exists}: the batch engine evaluates every
rate in milliseconds, the reference integrator produces trajectories, and the
outlier census (135 channels, all class-explained) already identifies which
channels are uncertain. A Monte Carlo over rate perturbations on a few hundred
states in the QSE window is a day's compute.

It is the cheapest of the three routes to answering \emph{``is $3\times10^{-6}$
the right gate?''} --- the others being \reg{T-01} (the output derivative,
partly infeasible per \S\ref{sec:explodability}) and \foreignreg{2}{R-11} (the
split-error comparison, which needs an experiment). The audit's recommendation
is to do this one first and let it inform the other two.

% ---------------------------------------------------------------------
\subsubsection{Axes checked and found already covered}
\label{sec:axescovered}

Listed so the sweep is auditable rather than merely asserted:

\begingroup
\small
\begin{tabularx}{\textwidth}{@{}L{6.4cm} Y@{}}
\toprule
\textbf{axis} & \textbf{covered by} \\
\midrule
plasma EOS, degeneracy, Coulomb coupling & Tier~0 \S 0.1 \\
statistical mechanics, detailed balance & Tier~0 \S 0.2, Tier~1 \S III \\
rate theory, Gamow, Hauser--Feshbach & Tier~1 \S I \\
screening & Tier~1 \S IV, with the two-$\kappa$ rule as executable policy \\
the weak sector, EC in a degenerate plasma & Tier~1 \S V \\
stoichiometry, conservation, projection & Tier~2 \S I \\
cancellation, $\kappa$, entropy production & Tier~2 \S II \\
NSE/QSE algebra & Tier~2 \S III \\
stiff integration, the augmented state & Tier~2 \S IV \\
dataset identity, join keys, censoring & Tier~0 \S III, \S\ref{sec:shippedpipeline} here \\
CRN theory, flux cones, deficiency & Tier~2 \S I.2b \\
the speedup ceiling under fallback & Tier~2 \S 0.5b \\
$e^{+}e^{-}$ pairs at the hot corner & Tier~1 \S VIII.C.6 (promoted to a decision) \\
isomers, ground-state vs stellar rates & Tier~1 \S VIII.C.2, \S VIII.C.8 \\
neutrino loss channels & Tier~0 \S 0.4, Tier~1 \S V.8 \\
learning theory / operator learning & \foreignreg{2}{R-24} (sketch) \\
rollout accumulation theory & Tier~0 \S IV.3 \\
GNN architecture specifics & Tier~4 sketch, \tierone{VIII.E.1} \\
\bottomrule
\end{tabularx}
\endgroup

Eighteen axes, all genuinely covered. The audit's yield (ten new areas) came
entirely from asking about the \emph{boundary} of the physics rather than its
interior.

% ---------------------------------------------------------------------
\subsubsection{Is the audit converging?}
\label{sec:converging}

Honestly: \textbf{the physics audit has converged; the boundary audit has not.}

Evidence for the first: passes A--C found exactly one physics-internal gap
(mixing), and it was already flagged as a caveat in Tier~0 --- the audit
upgraded a known approximation to a quantified one rather than discovering
something unknown. Four tiers of derivation have found essentially everything a
physicist would ask.

Evidence against the second: pass B (the interface) found four items from four
questions, at a 100\% hit rate, and pass C found four more. \textbf{A pass that
finds something every time it looks has not converged}, and the reason is
structural: the project has no artifact yet, so there has been nothing to ask
interface questions \emph{of}. Those questions will keep producing answers until
an emulator exists and is called from something.

The practical reading: \textbf{expect another audit's worth of items when the
first model is deployed, and budget for it.} The right response is not to try
to enumerate them now --- that is how you get a speculative list --- but to
notice that the register's cheapest and highest-ranked items
(\S\ref{sec:ranked}; \reg{T-24}, \reg{T-03}, \reg{T-08}) are all \emph{scoping}
decisions that would make the eventual interface work smaller.
% [edition] reworded cross-ref: "the current register's cheapest and highest-ranked items (**T-24**, **T-03**, **T-08**)" -- the register follows this section in the report
Doing scope before capability is the correct order and the register already
reflects it.

% =====================================================================
\subsection{The third pass: the emulator as an object}
\label{sec:emulatorobject}

\S\ref{sec:converging} predicted that another audit's worth of items would
appear once there was an artifact to ask questions of, and that the interface
axis had not converged. This part is that prediction tested early --- a third
pass run \emph{without} waiting for a model, by asking what one would measure
\textbf{of} a trained emulator and what a stellar code would do \textbf{with}
it.

It found ten more. They are not scattered: \textbf{nine of the ten are about the
model as a dynamical object or about its interaction with a host solver}, which
is the same region pass B found. That localisation is the useful result ---
\S\ref{sec:convergingnow}.

% ---------------------------------------------------------------------
\subsubsection{The emulator has its own fixed points, and nobody has asked what they are}
\label{sec:fixedpoints}

This is the strongest item in the pass, and it is a direct generalisation of the
tier's own headline finding.

\textbf{The true map's fixed-point structure.} At constant $(T,\rho)$ the
one-step map $\Phi_{\Delta t}$ is the time-$\Delta t$ flow of an autonomous ODE,
and its fixed points are that ODE's equilibria. For the strong/EM sector that is
\textbf{NSE at the current \Ye} --- a one-parameter family, since the strong
sector cannot move \Ye{} (\S\ref{sec:yeidentity}; the \Ye-parametrised NSE of
\citealt{hix1996} --- physics reasoning here, not a quoted result). Include the
weak sector and the family collapses to a single point ($\beta$-equilibrium),
but only on timescales far longer than anything here. So:
%
\begin{equation}
\label{eq:nseprojection}
\Phi_{\Delta t}(X) \;\xrightarrow[\ \Delta t \gg \tau_{\text{strong}}\ ]{}\;
X_{\mathrm{NSE}}\big(T,\rho,Y_e(X)\big) ,
\end{equation}
%
i.e.\ \textbf{at large $\Delta t$ the true map is very nearly a projection onto
a one-dimensional attracting manifold indexed by \Ye.}

\textbf{The learned map's fixed-point structure is whatever the weights make
it.} $\Phihat_{\Delta t}$ is some smooth function; its fixed set
$\{X : \Phihat(X) = X\}$ is determined by parameters fitted to minimise a
per-step loss. Nothing in the training objective mentions fixed points. There is
no reason for the learned attractor to be the NSE manifold, to be
one-dimensional, to be indexed by \Ye, or to exist at all.

Four consequences, in increasing order of how much they should worry someone:

\begin{enumerate}
\item \textbf{A long rollout converges to the learned attractor, not the
  physical one.} Per-step accuracy controls the transient; the fixed-point
  structure controls the destination. For $N \approx 10^{3}$ steps in a strongly
  contracting system, the destination is most of the answer.
\item \textbf{Per-step loss is nearly blind to it.} Near an attracting fixed
  point the true increment is \emph{small}, so a model can have tiny per-step
  error there and still put its own fixed point somewhere else entirely --- the
  errors that relocate an attractor are exactly the ones a small-increment loss
  weighs least.
\item \textbf{This is the same failure the labels have.} \S\ref{part:pathology}
  found the \emph{label generator} sitting at a displaced fixed point and called
  it a pathology. A trained emulator can acquire a displaced fixed point of its
  own, from a different cause, with the same consequence --- and it would be
  invisible to every gate the project currently specifies.
\item \textbf{\Ye{} is the parameter of the true attractor manifold.} If the
  learned attractor does not depend on \Ye{} in the right way, the model gets
  the \emph{endpoint} of silicon burning wrong as a function of the very
  quantity the project exists to predict (\S\ref{sec:endpoint}).
\end{enumerate}

\textbf{And it is trivially measurable, with tools already built.} Iterate the
trained map to convergence at fixed $(T,\rho)$ from many starting compositions,
and compare the limit against \code{qse.solve\_nse}:
%
\begin{equation}
\label{eq:fixedpointcensus}
X_\infty \;=\; \lim_{k\to\infty}\hat\Phi_{\Delta t}^{\,k}(X_0)
\quad\text{(iterate until } \|\Delta X\| < \mathrm{tol}\text{)},
\qquad
d \;=\; \big\|\log X_\infty - \log X_{\mathrm{NSE}}\big(T,\rho,Y_e(X_\infty)\big)\big\|_\infty ,
\end{equation}
%
sweeping $X_0$ over many starting compositions to find how many distinct
attractors exist and how far each sits from the Saha solution.

That is \code{scripts/step6\_label\_nse\_census.py} pointed at a model instead
of at a dataset --- the same census, the same reference solver, the same
statistics. It is also the direct extension of \reg{T-29}'s fixed-point audit
from datasets to models, and it should be a standing gate, not a one-off.

\textbf{It also suggests a training term nobody has proposed.} Add a fixed-point
consistency penalty $\|\Phihat(X_{\mathrm{NSE}}) - X_{\mathrm{NSE}}\|$
evaluated at Saha solutions --- free supervision (the Saha solver is cheap and
exact), aimed precisely at the property the per-step loss cannot see. It pins
the model's attractor to the physical one without touching the conservation map.

Note how this sits against \foreignreg{2}{R-13}: a second-law-consistent
emulator cannot \emph{diverge} --- the Helmholtz free energy bounds the fast
sector --- but a Lyapunov bound says nothing about \textbf{which} fixed point
the descent reaches. R-13 gives stability; this gives correctness of the
destination. They are complementary and only the first is on the register.
Registered as \reg{T-33}.

% ---------------------------------------------------------------------
\subsubsection{The learned Jacobian --- does stiffness survive learning?}
\label{sec:jacobian}

The companion diagnostic, equally cheap and equally absent.

\textbf{What the true map's Jacobian looks like.} If $J = \nu\,\partial R/\partial Y$
is the ODE Jacobian with eigenvalues $\lambda_i$ (Tier~0 \S V.3: spread
$10^{12}$--$10^{15}$), then the flow map's Jacobian has eigenvalues
$\approx e^{\lambda_i \Delta t}$ (the standard linearisation of a stiff flow,
\citealt{hairer1996}). With most $\lambda_i$ strongly negative and $\Delta t$
large:
%
\begin{equation}
\label{eq:flowspectrum}
\text{spec}\!\left(\frac{\partial \Phi_{\Delta t}}{\partial X}\right)
\;\approx\; \big\{\,e^{\lambda_i \Delta t}\,\big\}
\;\longrightarrow\; \{\underbrace{\approx 0,\dots,\approx 0}_{\text{the stiff directions}},\ \underbrace{O(1)}_{\text{the slow manifold}}\} .
\end{equation}

\textbf{The true one-step map is severely rank-deficient}: it annihilates almost
every direction and preserves a handful. That is the same statement as
``stiff'', as ``QSE'', and as ``the reachable set has effective dimension
3.4 / 8.6'' (\foreignreg{2}{R-02}, Tier~2 \S I.10b \derivedthere) --- this
tier's \S\ref{sec:mechanism} list of four names for one mechanism, now seen
through the flow map's spectrum.

\textbf{What that demands of the emulator, as two failure modes with opposite
signs:}

\begingroup
\small
\begin{tabularx}{\textwidth}{@{}L{3.6cm} L{3.6cm} Y@{}}
\toprule
\textbf{failure} & \textbf{signature} & \textbf{consequence} \\
\midrule
\textbf{under-contraction} --- learned singular values $> 1$ where the true map
  has $\approx 0$ &
  the model preserves directions physics destroys &
  rollout amplifies its own error along stiff directions: the classic
  autoregressive blow-up \citep{brandstetter2022}, and precisely what S20's
  pushforward loss \citep{brandstetter2022} and noise injection
  \citep{sanchezgonzalez2020} are \emph{heuristics} for \\
\textbf{over-contraction} --- learned singular values $\approx 0$ where the true
  map has $O(1)$ &
  the model collapses the slow manifold too &
  it destroys the \Ye{} signal and the bridge flow --- the physics --- while
  looking beautifully stable \\
\bottomrule
\end{tabularx}
\endgroup

Both are measurable in one line with autodiff, against a reference the project
already produces: \code{fluxes/integrate.py} can supply the true
$\partial\Phi/\partial X$ by forward sensitivity, and its analytic
$\partial R/\partial Y$ is already implemented (\code{ADR 0006}).

\textbf{Why this matters more than a diagnostic usually would.} S20 is currently
specified entirely in terms of \emph{training tricks} --- pushforward losses,
noise injection restricted to non-equilibrated channels, the Ono \& Sugimura
governor --- with no measurement that says whether the trained map is
contracting on the right subspace. A spectral comparison turns rollout stability
from a hoped-for property into a checked one, and it does so \textbf{before} an
expensive rollout study rather than after. Registered as \reg{T-34}.

% ---------------------------------------------------------------------
\subsubsection{The fallback gate is a discontinuity in the host's implicit solve}
\label{sec:dispatcher}

A specific and serious interaction that no document anticipates.

MESA's structure solve is a \textbf{Newton iteration} over the whole stellar
model (\citealt[\S 6]{paxton2011}; \citealt{jermyn2023} --- the partial
derivatives are now supplied by automatic differentiation), and the network's
output enters the residual (through \enuc, the composition, and the EOS's
dependence on $\bar A$ and $\bar Z$). Newton needs a residual function that is
smooth in the iterates.

Now insert S22's design: an OOD detector computes a score $s(X, T, \rho)$ and
\textbf{dispatches} to either the emulator or the real solver. As Newton adjusts
the state during its iteration, $s$ changes; if it crosses the dispatch
threshold, the residual function \textbf{jumps} by the emulator-vs-solver
difference --- which is precisely the quantity the gate exists because it is
large.

The consequences are the standard ones for a discontinuous residual, and they
are not subtle: Newton fails to converge, or cycles between two branches, or
converges to a spurious point on one side. The host responds by \textbf{cutting
the timestep}, which changes the state, which may flip the dispatch again. This
is a well-known pathology with discontinuous EOS tables and phase transitions
(MESA blends its EOS and opacity tables, \citealt{paxton2011}, and had to
handle the discontinuous derivatives of the Skye crystallisation transition
explicitly, \citealt{jermyn2023}), and it is exactly reproduced here.

\textbf{Three mitigations, none currently specified:}

\begingroup
\small
\begin{tabularx}{\textwidth}{@{}L{7.2cm} Y@{}}
\toprule
\textbf{mitigation} & \textbf{cost} \\
\midrule
\textbf{freeze the dispatch for the duration of one Newton solve} (decide once
  from the incoming state) &
  free; probably the right default \\
\textbf{blend continuously in a band} --- a weighted average of emulator and
  solver over a threshold interval &
  pays the solver's cost across the whole band, worsening the effective $f$ of
  the $1/f$ ceiling \\
\textbf{require agreement to within solver tolerance in the switching band}
  --- i.e.\ only switch where it does not matter &
  strongest, but constrains where the gate is allowed to fire \\
\bottomrule
\end{tabularx}
\endgroup

\textbf{Credit where it is due, and the distinction.}
\code{docs/reports/consolidated-gnn-architecture.md} \S 5 already identifies
that \emph{mask membership flips} inject discontinuities in $\dd Y$ and \enuc{}
that destabilise \textbf{the rollout gradient}, and proposes churn penalties,
annealed gates and hysteresis. That is the same mathematics one level up,
applied to training. What is new here is that the same problem appears
\textbf{at the host's Newton solve at deployment}, where the switch is not the
mask but the \emph{fallback dispatcher}, and where the remedies available during
training (annealing, soft gates) are unavailable. Registered as \reg{T-35}.

% ---------------------------------------------------------------------
\subsubsection{\texorpdfstring{$\rho$}{rho} is not a free input dimension}
\label{sec:rhoinput}

An architecture insight that falls out of the rate structure and that nothing in
the project uses.

Write the ODE by reaction order. At fixed $T$, with $Y$-dependence factored out:

\begingroup
\small
\begin{tabularx}{\textwidth}{@{}L{5.2cm} Y L{2.6cm}@{}}
\toprule
\textbf{sector} & \textbf{contribution to $\dd Y/\dd t$} & \textbf{$\rho$-dependence} \\
\midrule
one-body (photodisintegration, $\beta$-decay) & $\lambda(T)\cdot Y$ & \textbf{none} \\
two-body (captures) & $\rho\, N_A\langle\sigma v\rangle(T)\cdot Y_aY_b$ & \textbf{$\propto \rho$} \\
three-body (triple-$\alpha$ and friends) & $\rho^{2}\ldots\cdot Y_aY_bY_c$ & \textbf{$\propto \rho^{2}$} \\
tabulated weak (EC) & $\lambda(T, \rho\Ye)\cdot Y$ & through \textbf{$\rho\Ye$} \\
\bottomrule
\end{tabularx}
\endgroup

Change variables to $\tau = \rho t$. Then the two-body sector becomes
$\rho$-free, the three-body sector carries one power of $\rho$, the one-body
sector carries $\rho^{-1}$, and the tabulated sector carries whatever the table
does. So:

\begin{result}
\textbf{If the network contained only two-body reactions, the map would be
exactly invariant under $(\rho, \Delta t) \to (\lambda\rho, \Delta t/\lambda)$.}
$\rho$ would not be an independent input at all --- only the product
$\rho\Delta t$ would matter. (With screening off: the screening enhancement
depends on $\rho$ through $\Gamma$ and breaks the invariance.)
\end{result}

The real network is not two-body-only, and the symmetry-breaking terms are
\emph{exactly the physically interesting ones} --- photodisintegration is what
creates the equilibrium (\S\ref{sec:ignition}), and electron capture is what
moves \Ye. That is the point rather than an objection, and it yields three
usable consequences:

\begin{enumerate}
\item \textbf{A better conditioning variable.} The model is currently
  conditioned on $\log\rho$ and $\log\Delta t$ as separate features.
  \textbf{$\rho\Delta t$} captures the entire two-body sector --- the numerical
  majority of columns (ch.~4--5 alone are 397/979 of the reactions
  \resultsref{2026-07-08}; ch.~6--7 are two-body too) --- in one variable,
  leaving the network to learn only the \emph{departure} from that scaling.
  This composes with \foreignreg{2}{R-10}'s $\Delta t/\tau$ proposal rather
  than competing with it: R-10 says condition on how many relaxation times the
  step covers, and this says the density enters that relaxation through a known
  power law.
\item \textbf{An exact consistency test.} Evaluate the trained model at
  $(\rho, \Delta t)$ and at $(\lambda\rho, \Delta t/\lambda)$ and compare
  against the same pair evaluated on the reference integrator. The
  \emph{reference} violates the invariance by a known, computable amount (the
  one- and three-body contributions); the model's violation should match it. A
  model whose violation is the wrong size has learned the wrong decomposition,
  and this test needs no labels.
\item \textbf{A physical diagnostic for free.} The measured degree of invariance
  violation, per stratum, \textbf{is} a measure of how much of the local
  dynamics is photodisintegration-driven --- i.e.\ an independent, rate-free
  probe of exactly the QSE onset that \S\ref{sec:ignition} derives from $Q/kT$.
\end{enumerate}

Registered as \reg{T-36}.

% ---------------------------------------------------------------------
\subsubsection{What ``calibrated'' means for a deterministic map}
\label{sec:calibrated}

\reg{T-26} concluded that S22 needs a \emph{calibrated continuous error
estimate} rather than a binary flag. That conclusion has a literature and an
assumption problem, neither of which is in any document.

\textbf{The specified approach is deep ensembles}, and it has two independent
difficulties here:

\begin{itemize}
\item \textbf{Economics.} An ensemble of size $k$ costs $k\times$ per call, and
  every call --- not just the fallbacks. Tier~2 \S 0.5b's cost model puts
  $t_{\mathrm{det}}$ on the critical path and shows a detector costing 10\% of
  the solver caps the speedup at $10\times$ before a single fallback happens.
  An ensemble is the most expensive UQ there is, in a setting where
  \foreignreg{2}{R-12} shows the detector's cost is a first-order term.
\item \textbf{Statistics.} All UQ methods, deep ensembles
  \citep{lakshminarayanan2017} included, \textbf{lose calibration under
  distribution shift} --- ensembles degrade least in the benchmark of
  \citet{ovadia2019} but still substantially --- and distribution shift is not
  a risk here --- it is a \emph{measured fact} (\S\ref{sec:twodistributions}).
  (The earlier wording ``worst-calibrated under shift'' inverted the source;
  corrected in the 2026-08-18 audit.) The consolidated architecture report
  already names this as a load-bearing risk (``a jointly-overconfident ensemble
  could fail the OOD gate silently exactly where the transfer fails''), which
  is the right worry and has no proposed remedy.
\end{itemize}

\textbf{The tool the project has not named is conformal prediction}
\citep{vovk2005, angelopoulos2021}, which gives distribution-free,
finite-sample-valid prediction intervals from a calibration set. Its assumption
is exchangeability, and here the assumption fails in two \emph{known} ways,
which is the useful part:

\begingroup
\small
\begin{tabularx}{\textwidth}{@{}L{3.6cm} L{3.4cm} Y@{}}
\toprule
\textbf{failure} & \textbf{already measured as} & \textbf{the conformal variant that handles it} \\
\midrule
rows within a trajectory are correlated &
  ICC 0.22--0.70 (\reg{T-05}) &
  block / cluster conformal --- calibrate over \textbf{trajectories}
  \assumednote{e.g.\ the block-conformal methods of \citealt{chernozhukov2018},
  not opened here} \\
the deployment measure $\ne$ the training measure &
  \S\ref{sec:twodistributions}, \foreignreg{2}{R-07}, \reg{T-13} &
  \textbf{weighted} conformal \citep{tibshirani2019}, which needs the
  likelihood ratio between the two measures --- its covariate-shift form
  assumes $P(Y|X)$ unchanged, true for a deterministic emulator target \\
\bottomrule
\end{tabularx}
\endgroup

The second row is the striking one: \textbf{weighted conformal prediction
requires exactly the density ratio that R-07 / T-13 set out to measure.} A
sampling-measure study filed under ``is the Sobol grid representative?'' turns
out to be the input to a principled UQ method. Two open items that looked
unrelated are the same measurement.

\textbf{And there is a cheap alternative the project owns and has not connected
to S22.} It already computes several \emph{physics residuals} that require no
ensemble and cost essentially nothing:

\begin{itemize}
\item the \textbf{entropy-production sign} check (\foreignreg{2}{R-05}: a
  learned $\hat\varphi$ can violate the second law and pass every current
  gate);
\item the \textbf{energy-consistency} residual between the flux route and the
  composition route;
\item \textbf{positivity} violations (\foreignreg{2}{R-04});
\item the \textbf{fixed-point distance} of \S\ref{sec:fixedpoints}.
\end{itemize}

Each is a scalar, deterministic, single-forward-pass validity score derived from
physics rather than from statistics --- which is precisely the ``cheap UQ'' that
Tier~2 \S 0.5b argues is structurally preferred on economic grounds. \textbf{The
project has been building OOD scores for two tiers without labelling them as
such.} Registered as \reg{T-37}, and it is a reframing rather than a new
measurement.

% ---------------------------------------------------------------------
\subsubsection{Rates as inputs --- the capability nobody claimed}
\label{sec:ratesinputs}

A design question with a scientific payoff, unasked.

The emulator's inputs are $(T, \rho, X)$. The \emph{rates} are baked into the
weights. Suppose instead the model took a low-dimensional rate-perturbation
vector $\theta$ --- multipliers on a chosen set of channels, or a categorical
selector over weak-rate library families (FFN / Oda / LMP). Then autodiff gives
%
\begin{equation}
\label{eq:ratederiv}
\frac{\partial X'}{\partial \theta}, \qquad \frac{\partial Y_e'}{\partial \theta}
\end{equation}
%
\textbf{for free, at inference time.} That is exactly \S\ref{sec:nucleardata}'s
rate-uncertainty $\to$ \Ye{} propagation --- currently a day of compute per
study --- reduced to a milliseconds-per-state gradient evaluation, and it is the
kind of thing surrogate models are uniquely good for.

\textbf{Honest assessment of the cost.} Training data would have to \emph{vary}
the rates, which means regenerating labels under perturbed rate sets. For a full
continuous $\theta$ that is prohibitive. But two restricted versions are
affordable and useful:

\begin{itemize}
\item \textbf{A categorical library selector} with 2--3 values (the FFN/LMP pair
  being the obvious one, since that shift \emph{is} the project's accuracy
  floor, \S\ref{sec:errorhierarchy} row 4). This costs one extra label campaign
  per library and would let a user ask ``how much of my \Ye{} is the rate
  library?'' directly.
\item \textbf{A handful of scalar multipliers on the top-20 \Ye{} carriers}
  (\S\ref{sec:yecarriers}), which are few, identified, and dominate the budget.
\end{itemize}

Whether to do this in Phase 1 is a scope call and probably ``no''. Whether to
\emph{say} it is a scope call is not --- because it changes what the artifact is
\textbf{for}: from ``a fast network'' to ``a differentiable model of the
network's dependence on its own inputs''. That is a stronger scientific product
and it is not in any plan. Registered as \reg{T-38}.

% ---------------------------------------------------------------------
\subsubsection{Which \texorpdfstring{$\Delta t$}{dt} does the host actually use?}
\label{sec:hostdt}

The emulator is trained on nine values spanning $10^{-6}\ldots10^{2}$\,s,
inherited from the NNN's grid, which its authors chose to bracket the timesteps
of their three MESA test-case tracks (\citealt[\S 2.2 and Fig.~1]{grichener2025})
and extended ``several orders of magnitude'' below them ``which might be needed
for hydrodynamical simulations''. \textbf{Nobody has measured the distribution
of $\Delta t$ that MESA actually takes during silicon burning} --- the paper's
Fig.~1 is three tracks, not a per-phase histogram.

If the host spends most of its steps in $10^{-2}\ldots10^{0}$\,s, then three of
the nine decades are training capacity spent on regions never visited, and the
extrapolation risk lives at a different end of the grid than assumed. If the
host takes steps \emph{shorter} than $10^{-6}$\,s --- the NNN authors say the
grid's low end already lies below MESA's steps, so this would need a
hydrodynamic phase --- the emulator is being extrapolated from its first step,
and \S\ref{sec:realtrajectory}'s finding that even the shortest label is a full
stiff relaxation acquires a sharper edge.

\textbf{This is the same MESA run as \reg{T-24}}, so it is free once that
profiling exists: histogram \code{dt} by evolutionary phase alongside the
\code{net} timing. Registered as \reg{T-39}, and it should be done in the same
afternoon.

% ---------------------------------------------------------------------
\subsubsection{Sample complexity --- a million points in 82 dimensions}
\label{sec:samplecomplexity}

A framing that the project's own measurements make sharp, and that nothing
states.

The training corpus is ${\sim}10^{6}$ states per network, with inputs of
dimension $n + 2 = 82$ / 153. The naive covering estimate for a smooth function
on a $d$-dimensional set at resolution $\varepsilon$ is $\varepsilon^{-d}$
samples; at $d = 82$ the number is not worth writing down. \textbf{Uniform
coverage is not merely infeasible, it is absurd by seventy orders of magnitude.}

The reason it can nonetheless work is that the \emph{deployment} set is much
smaller than the input space: \foreignreg{2}{R-02} measures the effective
dimension of visited compositions at 3.4 / 8.6. At $d = 8.6$ and
$\varepsilon = 0.1$ the covering estimate is ${\sim}10^{8.6} \approx 4\times10^{8}$
--- still \textbf{two to three orders above the corpus}, though this is a
worst-case bound that smoothness and the map's near-projection structure
(\S\ref{sec:jacobian}) both improve substantially.

But the real point is a \emph{capacity allocation} argument, distinct from the
distribution-shift argument of \S\ref{sec:twodistributions} and, as far as these
notes can tell, not made anywhere:

\begin{aside}
The training compositions are \textbf{random}, hence spread over the full
high-dimensional simplex. The deployment compositions lie on a ${\sim}3$--9
dimensional manifold. So the model spends the overwhelming majority of its
capacity, and of the corpus, learning behaviour on states it will never be asked
about --- while the manifold it \emph{will} be asked about is covered by
whatever small fraction of the corpus happens to lie near it.
\end{aside}

That is an argument for trajectory-derived training data on grounds of
\textbf{sample efficiency}, independent of whether the distributions match. It
also predicts that a modest number of on-manifold states could be worth a great
many off-manifold ones --- which, if true, makes the $10^{3}$--$10^{4}$-state
$\Phifl$-label budget of \S\ref{sec:phisupervision} look much less like a
compromise and much more like the right design. Registered as \reg{T-40}.

% ---------------------------------------------------------------------
\subsubsection{Two smaller ones}
\label{sec:twosmaller}

\textbf{Network switching mid-run.} MESA can change nets during an evolution
(\code{star\_job} \code{change\_net} / \code{new\_net\_name}, r23.05.1
\code{star\_job.defaults:2364}; a star does not need an iron-peak network during
hydrogen burning). An emulator trained on \msn{80} has nothing to say about the
step at which the host switches nets, and the composition must be mapped between
species lists at that point. Minor, but it is a real code path and it interacts
with the size-transfer story. Folded into \reg{T-32}.

\textbf{Licensing, attribution, and data reuse.} The project redistributes
nothing but depends on Zenodo 14873443 (someone else's data and trained models),
MESA, bbq and pynucastro, and will publish a head-to-head comparison against the
NNN using the NNN authors' own shipped artifacts. The licence terms, the citation
obligations, and the norms around benchmarking someone's published model with
their own weights are a genuine publication prerequisite that appears in no
document. It is half an hour of checking and it is the kind of thing that is
embarrassing rather than fatal --- which is exactly why it gets skipped. Folded
into \reg{T-32}.

% ---------------------------------------------------------------------
\subsubsection{One unretired premise, found while sweeping}
\label{sec:unretired}

Not an area, but a finding, and it is the same \emph{class} as the ``6--8
orders'' timescale figure that \S\ref{sec:gates378} retired.

\code{docs/reports/consolidated-gnn-architecture.md} \S 1 states, as inherited
design basis:

\begin{quote}
``the net flow between the silicon and iron-peak groups is carried predominantly
by the single bottleneck reaction \textbf{\nuc{Sc}{45}(p,$\gamma$)\nuc{Ti}{46}}
(\textbf{$\approx$75\% of the flow} at the proton-rich $Z = 21$ edge), which is
therefore the reaction whose flux the Phase-0 kill-test instruments.''
\end{quote}

The measurement \resultsref{2026-07-12} says: \nuc{Sc}{45}(p,$\gamma$)\nuc{Ti}{46}
is \textbf{rank 2 of 251} inter-group carriers (\msn{151}) with \textbf{share
0.10--0.12} under the \code{a24\_46} boundary on the full relaxed manifold
(0.094 on the 2026-07-10 pre-stall high-\Ye{} QSE-window subset), and the
\textbf{top-20 columns carry 0.88} (the \msn{80} bridge figure,
\resultsref{2026-07-12}).

Both the verdict document and \S\ref{sec:twoclusters} of this file report that
as \textbf{CONFIRMED} --- and it is, as a confirmation that the reaction is a
leading carrier, which is a genuine and pleasing independent recovery of a
literature claim. But \textbf{the magnitude the design leaned on is off by a
factor of six or seven}, and nothing retires it. ``Predominantly carried by a
single reaction'' and ``top-20 carry 0.88, of which the leader has 0.11'' are
different physical pictures, and they imply different architectures: the first
justifies a bottleneck-focused design, the second justifies exactly the
concentrated-but-not-singular loss weighting the project actually adopted.

Three things should happen: pin the 75\% figure to what it actually measures ---
it \emph{is} sourced (\citealt{woosley1973}, via \citealt{hix1996}: WAC
``contended that this linkage was dominated by a single reaction,
\nuc{Sc}{45}(p,$\gamma$)\nuc{Ti}{46}, with perhaps a quarter of the flow going
through less important reactions''; \citeauthor{hix1996}'s own value is
$\approx$70\% of the flux \emph{into the iron-peak group} at $\Tn = 5$,
$\rho = 10^{7}$, $\Ye = 0.498$, integrated while $X$(Si group) falls
$0.95 \to 0.85$), but it is a different quantity at a different state from
``share of all inter-group carriers on the relaxed manifold under
\code{a24\_46}''; retire the transplanted magnitude in \code{RESULTS.md} with
the measured replacement, exactly as the timescale figure was retired; and
audit the remaining inherited design premises for others of the same kind. The
two found so far share a pattern worth naming: \textbf{a magnitude attached to a
correct qualitative claim, either unsourced (the 6--8 orders) or sourced for a
different quantity and state (the 75\%)} (corrected in the 2026-08-18 audit).
Registered as \reg{T-41}, and it is a \code{referee}-subagent task.

% ---------------------------------------------------------------------
\subsubsection{Is the audit converging \emph{now}?}
\label{sec:convergingnow}

Better than after pass 3, and the reason is the shape of the yield rather than
its size.

\textbf{Pass 3} (\S\ref{sec:boundaryaudit}) found ten areas spread across four
axes: physics-adjacent (mixing), setting (progenitors, NSE handoff,
explodability), interface (four items), and practice (three items). Spread
means \emph{not localised} means \emph{not converging}.

\textbf{This pass} found ten items of which \textbf{nine sit in one place}: the
emulator considered as a dynamical system (\S\ref{sec:fixedpoints},
\ref{sec:jacobian}, \ref{sec:rhoinput}, \ref{sec:samplecomplexity}), and its
coupling to a host solver (\S\ref{sec:dispatcher}, \ref{sec:calibrated},
\ref{sec:ratesinputs}, \ref{sec:hostdt}). The tenth (\S\ref{sec:unretired}) is
a documentation finding, not an area. That is what localisation looks like, and
it supports a specific prediction:

\begin{result}
\textbf{The remaining unexplored prerequisite is essentially one subject: the
theory and practice of a learned flow map --- its fixed points, its spectrum,
its conditioning symmetries, its calibration, and the numerics of embedding it
in an implicit solver.} That is a Tier-4 Part 0, and it is exactly the ``no
study node at all'' gap \tierone{VIII.E.1} identified from the architecture
side.
\end{result}

Two further observations that make the convergence claim honest rather than
convenient:

\begin{itemize}
\item \textbf{Diminishing novelty, not diminishing count.} This pass found as
  many items as the last, but they are far more homogeneous, and four of them
  (\S\ref{sec:fixedpoints}, \ref{sec:jacobian}, \ref{sec:calibrated},
  \ref{sec:samplecomplexity}) are \emph{reframings of things already measured}
  rather than new territory: the fixed-point census is \S\ref{part:pathology}'s
  method aimed at a model; the spectral test is stiffness seen through the flow
  map; conformal weighting needs R-07's density ratio; sample complexity needs
  R-02's effective dimension. \textbf{The project increasingly already owns the
  inputs to its own open questions.}
\item \textbf{The one thing that would genuinely change the picture is having
  an artifact.} Every item in \S\ref{sec:fixedpoints}--\ref{sec:calibrated} is
  a measurement \emph{of a trained model}, and none can be run before one
  exists. So the correct reading is not ``keep auditing'' but ``the audit has
  identified what to measure the moment there is something to measure, and that
  list should be a gate on the first training run rather than a
  retrospective.''
\end{itemize}

\textbf{The recommendation the two passes converge on}, stated once:
\emph{before the first training run, write the Tier-4 Part 0 and the
measurement plan it implies.} Sections \S\ref{sec:fixedpoints}--\ref{sec:calibrated}
and \foreignreg{2}{R-03}, \textbf{R-05}, \textbf{R-13}, \textbf{R-21}, \textbf{R-24} are
its contents; four of the eight
open Phase-0 checklist rows are its measurements; and two of the four
architecture components have retired premises waiting there
(\S\ref{sec:retired}).
% [edition] cross-ref "(SSVII.2)" -> \S\ref{sec:retired}, a forward reference in this report
